\silentchapter{5 Det finst agentar}{det finst agentar}
\prop{5}{}{Det finst agentar som responderer på landskapet\footnote{Ein agent kan vere alt frå ein ruter som handsamar trafikk, til ein brukar i eit grensesnitt.}.}
\prop{5.1}{a}{For kvar funksjonell klasse finst det minst eitt system som navigerer tilpassingslandskapet via negativ tilbakekopling\footnote{Eit kvart sjølv-justerande nettverk stettar krava til systemisk agens i landskapet.}.}
\prop{5.11}{d}{Ein \textbf{agent\footnote{Avstandsmålinga i digitale system er ofte eksplisitt definert via feilratar og responstider.}} er ein operator definert uttømande ved trippelet: måltilstand, avstandsmåling, justeringsmekanisme.\footnote{Rosenblueth, Wiener \& Bigelow (1943)}}
\prop{5.111}{t}{Definisjonen krev tilbakekopling.}
\prop{5.1111}{t}{Ho krev korkje medvit, intensjon eller nervesystem; ho krev berre funksjonell organisering.\footnote{Wiener (1948); Turing (1950)\footnote{At agens er substratuavhengig, er aksiomet som gjer at vi kan overføre formalisme frå biologi direkte til tenestedesign.}}}
\prop{5.11111}{t}{Eit system som ikkje treng å vite kva det er laga av for å navigere mot eit mål, er ein agent uavhengig av kva det er laga av.}
\prop{5.11112}{i}{Ein stein som rullar nedover ei skråning er ikkje ein agent.}
\prop{5.111121}{i}{Rørsla hans responderer ikkje på nokon avstand til nokon måltilstand.}
\prop{5.111122}{t}{Grensa mellom agent og ikkje-agent går ved informasjonsflyt.}
\prop{5.12}{o}{Agentar skil seg i tre målbare eigenskapar.}
\prop{5.121}{o}{Dei skil seg i kor mykje av formrommet dei kan representere, kor rask tilbakekoplinga er, og kor djup læringsevna er\footnote{Skilnaden mellom ein algoritme og eit operativsystem ligg ikkje i materie, men i spennvidda av deira kognitive lyskjegle og læringsdjupne.}.}
\prop{5.1211}{t}{Definisjonen er ikkje vakuøs; ho er stratifisert.}
\prop{5.12111}{t}{Differansen mellom ein termostat og ein meisterhåndverkar ligg ikkje i substrat, men i kompleksiteten i navigasjonskapasiteten.}
\prop{5.13}{t}{Agenten responderer ikkje på forma, men på landskapet.}
\prop{5.131}{t}{Måltilstanden er ikkje ein posisjon agenten finn opp\footnote{Routing-algoritmen finn ikkje opp den kortaste vegen; han oppdagar gradienten i nettverksforskingas formrom.}.}
\prop{5.1311}{t}{Det er ein haug i landskapet agenten oppdagar.}
\prop{5.13111}{t}{Agenten treng ikkje å forstå landskapet for å navigere det.}
\prop{5.131111}{t}{Det er tilstrekkeleg å registrere den lokale gradienten.}
\prop{5.131112}{t}{Navigasjonskompetanse krev ikkje kartografi.}
\prop{5.13112}{t}{To agentar produserer same resultat om dei registrerer same lokale gradient.}
\prop{5.131121}{t}{Dei treng ikkje ha nokon representasjon av det globale landskapet.}
\prop{5.131122}{t}{Kompetansen er lokal; effekten er global.}
\prop{5.2}{d}{Den \textbf{kognitive lyskjegla\footnote{For eit mikroteknestekomponent er lyskjegla strengt identisk med dei variablane som vert mottatt i HTTP-forespurnaden.}} til ein agent er den delmengda av formrommet han kan representere og påverke.\footnote{Fields \& Levin (2022)}}
\prop{5.21}{t}{Det som ligg utanfor lyskjegla, er kausalt utilgjengeleg\footnote{Alt utanfor dette datalaget er, for komponenten, kausalt eit svart hòl.}.}
\prop{5.211}{o}{Lyskjeglene varierer i skala.}
\prop{5.2111}{o}{Dei strekkjer seg frå cellulær fysiologi over minutt til designtradisjonar over hundreår.}
\prop{5.22}{t}{Lyskjegla avgrensar oppløysinga.}
\prop{5.221}{t}{Kva delformer ein agent kan identifisere, avgjer kva transformasjonar som er tilgjengelege\footnote{Grensesnittets oppløysing bestemmer kva operasjonar brukaren kan ta: grovare UI, redusert kompetanse og færre transformasjonar.}.}
\prop{5.2211}{o}{Barnet, snikkaren og algoritmen dekomponerer same form i ulike operasjonelle einingar.}
\prop{5.22111}{o}{Same form; ulik oppløysing; ulike tilgjengelege operasjonar.}
\prop{5.2212}{t}{Ein agent oppdagar ikkje former; han oppdagar affordansar.}
\prop{5.22121}{t}{Ein affordans er ein lokal eigenskap ved landskapet som berre er synleg innanfor lyskjegla\footnote{Affordansane i ei plattform spring ut av dei formelle kryssingane mellom brukar-lyskjegla og system-lyskjegla.}.}
\prop{5.22122}{t}{Å utvide lyskjegla er å gjere latente affordansar synlege.}
\prop{5.22123}{t}{Å krympe lyskjegla er å gjere realiserte affordansar usynlege.}
\prop{5.23}{d}{Fem operasjonar modifiserer lyskjegla: utviding, oppløysingsauke, rørsle, kollaps, skjerping\footnote{Å gje ein teneste nye tilgangsrettar er, matematisk sett, ein utviding av agensens lyskjegle.}.}
\prop{5.231}{d}{Utviding: lyskjegla dekkjer nye regionar. Handverkaren oppdagar eit nytt materiale.}
\prop{5.232}{d}{Oppløysingsauke: fleire delformer vert synlege i same region. Snikkaren lærer å sjå nye samanføyingar.}
\prop{5.233}{d}{Rørsle: lyskjegla flyttar seg til ukjent territorium. Designaren går frå møbel til arkitektur.}
\prop{5.234}{d}{Kollaps: lyskjegla krympar til éin realisert posisjon. Avgjersla vert endeleg.}
\prop{5.235}{d}{Skjerping: same dekning, høgare presisjon. Meistaren foredlar teknikken sin.}
\prop{5.2351}{t}{Operasjonane føreset berre kausal rekkevidd og tilbakekopling.}
\prop{5.23511}{t}{Dei krev ingen representasjon av heile formrommet.}
\prop{5.3}{d}{Blant agentens kompetansar er evna til å operere direkte på geometri\footnote{Formgrammatikken opererer her som ei rekkje reglar for API-anrop: identifiser del-tilstand, erstatt med oppdatert tilstand.}.}
\prop{5.31}{d}{Mengda av former i eit euklidsk rom, utstyrt med operasjonane sum og differanse, utgjer ein \textbf{formalgebra}.}
\prop{5.311}{t}{Formalgebraen er grunnlaget for det generative systemet.}
\prop{5.3111}{d}{Ein \textbf{formgrammatikk} er eit sett reglar som opererer direkte på geometri.}
\prop{5.31111}{d}{Han finn ei delform i den noverande forma og erstattar ho med ei anna.\footnote{Stiny \& Gips (1972)}}
\prop{5.3112}{t}{Kvar regel er definert under innleiring\footnote{Tenestelogikken si 'innleiring' krev berre at dei formelle parametrane er oppfylte for at operasjonen skal eksekverast; den historiske stien dit er sletta.}.}
\prop{5.31121}{t}{Ei delform kan identifiserast i den noverande forma uavhengig av korleis ho vart bygd.}
\prop{5.31122}{t}{Reglane ser; dei hugsar ikkje.}
\prop{5.31123}{t}{Forma sjølv avgjer kva reglar som kan fyre. Den noverande geometrien er den einaste premissen. Konstruksjonshistoria er irrelevant.}
\prop{5.32}{t}{Grammatikken genererer ein naboskap for kvar form.}
\prop{5.321}{d}{Naboskapen er mengda av alle former som kan nåast ved éin operasjon.}
\prop{5.3211}{d}{Han utgjer formens lokale moglegheitsrom.}
\prop{5.32111}{t}{Naboskapen er ikkje symmetrisk. At du kan kome dit, tyder ikkje at du kan kome attende.}
\prop{5.32112}{t}{Formrommet har ein retta struktur.}
\prop{5.32113}{t}{Kombinasjonen av to former kan framkalle delformer ingen av dei inneheldt åleine.}
\prop{5.321131}{t}{Same form, ulike dekomposisjonar. Dette er ikkje fordi forma er tvetydig, men fordi ho er rikare enn kvar einskild operasjon som produserte ho\footnote{Ein datasettkombinasjon kan avsløre mønster (delformer) som ingen av kjeldene bar i seg sjølve. Relasjonsdatabasens magi kviler på denne proposisjonen.}.}
\prop{5.33}{t}{Grammatikken spesifiserer kva som er mogleg.}
\prop{5.331}{t}{Seleksjonstrykka avgjer kva som vert realisert.}
\prop{5.3311}{t}{Grammatikken og landskapet møtest i agentens val.}
\prop{5.33111}{t}{Naboskapen gjev moglegheitene, landskapet gjev kvar moglegheit ein verdi.}
\prop{5.33112}{t}{Agenten aktualiserer den forma som stettar gradienten.}
\prop{5.33113}{t}{Sidan lyskjegla avgjer kva innleiringar agenten kan oppdage, er den generative kapasiteten avgrensa av lyskjegla.}
\prop{5.331131}{t}{To agentar med same grammatikk men ulike lyskjeglar produserer ulike former.}
\prop{5.4}{d}{Materialet er ein agent av null-orden\footnote{Maskinvaren og internettprotokollen fungerer som null-ordens agentar: dei tvingar all programvare til å unngå nettverksavbrot og minnelekkasjar.}.}
\prop{5.41}{t}{Det navigerer ikkje mot ein måltilstand.}
\prop{5.411}{t}{Det utelukkar posisjonar og favoriserer andre gjennom sine affordansar.}
\prop{5.4111}{t}{Materialet gjev landskapet sin topografi.}
\prop{5.41111}{i}{Ein planaria navigerer morforommet via bioelektriske gradientar.\footnote{Levin (2022, 2025)}}
\prop{5.41112}{i}{Ein handverkar navigerer formrommet via taktil tilbakemelding.}
\prop{5.41113}{i}{Ein diffusjonsmodell navigerer eit latent rom via denoising.}
\prop{5.41114}{t}{Substrata er ulike; strukturen er identisk.}
\prop{5.5}{t}{Navigasjonskompetanse er akkumulert.}
\prop{5.51}{t}{Kvar realisert form utvidar agentens lyskjegle.}
\prop{5.511}{d}{Læring er den temporale utvidinga av agentens tilgjengelege formrom.}
\prop{5.5111}{t}{Læring er hierarkisk: frå enkel navigasjon via strategiskifte til restrukturering av sjølve formrommet.}
\prop{5.51111}{t}{Læring, biologisk seleksjon og bayesiansk inferens er isomorfe operasjonar i ulike substrat\footnote{Nettverksruting, evolusjon og backpropagation-læring i nevrale nettverk er isomorfe; dei utfører den same algoritmiske avstandskrympinga i ulike substrat.}.}
\prop{5.51112}{t}{Dei minskar alle avstanden mellom modell og verd.}

